\chapter{Conclusion}
\label{chap:conclusion}

This dissertation investigated whether an intrusion detection system for RPL networks remains reliable when attack behaviour changes. The work moved beyond a fixed supervised benchmark by implementing nine attack families in Contiki-NG/Cooja, preserving matched attack and control runs, and evaluating transfer between distinct protocol mechanisms.

The baseline campaign produced 90 validated simulations and a frozen dataset of 810 labelled routing windows. The audit preserves run provenance, verifies the 240-second activation boundary, and excludes explicit implementation markers from model inputs. This dataset and its reproducible extraction pipeline form one practical contribution of the project.

The principal finding is that static IDS performance transfers poorly between RPL attack mechanisms. CART produced zero attack recall in 48 of 72 ordered cross-attack pairs, with mean F1 of 0.1822. Gaussian transfer was weaker. These failures show that strong performance within one attack environment does not demonstrate robustness to a changed attack surface.

Whole-seed adaptation substantially recovered performance. One complete target seed increased mean CART F1 to 0.8258, while three target seeds produced 0.8792. Because complete target seeds were reserved for evaluation, this improvement does not depend on random row-level leakage. The result supports adaptation but also shows its dependence on representative features and trustworthy labels.

Sybil identity manipulation strengthened the original contribution. It added a mechanism that changes apparent RPL control-plane identity rather than forwarding, rank, or route choice alone. Static Sybil-target performance was weak, while one complete Sybil adaptation seed produced mean held-out F1 of 0.9683. The claim remains bounded to the implemented Cooja scenario.

The online experiment showed that protocol-aware rank state can identify a focused blackhole-to-sinkhole change. CUSUM and delayed-error DDM detected all five evaluated sinkhole runs without matched-control alerts, with CUSUM responding one window earlier. The robustness campaign then showed that attacker position and radio loss can materially alter static behaviour, including severe false-positive failures. Target-condition adaptation recovered strong performance in the tested seeds.

The sinkhole defence experiment established a crucial distinction between detection and mitigation. Minimum-rank evidence was useful for monitoring, but direct parent avoidance reduced application responses by 79.06 percent and increased radio events by 571.22 percent. A credible defence requires temporal confidence, hysteresis, bounded route changes, recovery, and proper resource measurement.

Relative to SVELTE, this work retains the importance of protocol-aware routing evidence but focuses on cross-attack transfer and adaptation rather than a deployable mapper architecture. Relative to the Gope adversarial RL framework, it provides a simpler Cooja-based experiment with explicit attack code, activation timing, seed provenance, and leakage controls. It does not reproduce the full learned adversary or defender architecture.

The dissertation therefore contributes a controlled method for exposing static IDS brittleness, measuring seed-safe recovery, introducing an identity-related attack surface, and critically testing whether detector evidence can support mitigation. Future work should replicate the protocol on compatible external datasets, scale the Cooja topology and workload, test physical motes, evaluate mixed and mobile attacks, and design a stable trust-aware response. Any future LLM explanation layer should be assessed for evidence faithfulness and unsupported claims rather than treated as the detector.

The evidence also provides a practical demonstration of research iteration. Weak sinkhole features, invalid robustness builds, and harmful defence behaviour were not removed from the project history or converted into positive claims. They were used to improve instrumentation, isolate valid runs, and define limits. This critical treatment is central to the final contribution: the value of an adaptive IDS study lies not only in its highest F1 score, but in showing when a detector fails, what new evidence enables recovery, and why a seemingly reasonable response can make the network worse.
